\documentclass[10pt,a4paper]{amsart}
\usepackage[margin=19mm]{geometry}
\usepackage{amsmath,amssymb,mathtools}
\usepackage[scr=rsfs,cal=euler]{mathalfa}
\usepackage{xfrac}
\usepackage[unicode,hidelinks,bookmarks=false]{hyperref}
\newcommand{\ie}{i.e.\ }
\newcommand{\cf}{cf.\ }
\newcommand{\resp}{respectively\ }
\newcommand{\Subsection}{\subsection}
\providecommand{\nobreakdash}{\nobreak\hbox{-}}
\setlength{\emergencystretch}{2em}
\multlinegap0pt
\usepackage{amssymb}
\usepackage{mathtools}
\usepackage{enumerate}
\usepackage{tikz}
\usepackage{multirow}
\usepackage{xcolor}
\usepackage{booktabs}
\usepackage[all]{xy}
\newtheorem{proposition}{Proposition}
\newcommand{\Z}{\mathbb{Z}}
\newcommand{\op}[1]{\operatorname{#1}}
\newcommand{\tr}{\operatorname{tr}}
\title{Distribution questions for Isogeny Graphs over Finite Fields}
\author{Anwesh Ray}
\date{Source: arXiv:2512.14469v2 (CC BY 4.0)}
\begin{document}
\maketitle

\noindent\emph{Source and licence.} Distribution questions for Isogeny Graphs over Finite Fields. Version arXiv:2512.14469v2 (CC BY 4.0), distributed by arXiv under the Creative Commons Attribution 4.0 International licence, http://creativecommons.org/licenses/by/4.0/, as recorded in the arXiv licence metadata for this exact version and shown on the abstract page https://arxiv.org/abs/2512.14469v2. Full licence notice: \texttt{licenses/arxiv-2512.14469-CC-BY-4.0.txt}.\par\smallskip

\noindent\textit{Editorial scope.} Counted unit: the article's proposition prop (source0.tex lines 660--677). The article's complete original proof of it is delivered with it (source0.tex lines 679--808, one \texttt{proof} environment of 130 lines). No mathematics is written, repaired or re-derived: the statement and the proof are the article's own lines, in the article's own notation, and no retained line is substituted or re-flowed.  No further source-local context is needed: every cross-reference of every retained line resolves inside the two delivered spans.  Nothing was omitted from any retained span, so the package carries no omission record.  No retained line contains a citation, so the wrapper carries no bibliography and no bibliography entry.  No retained line contains a figure environment, an image-inclusion command or drawing code, so no image is delivered and the package declares no asset dependency.  Editorial formatting only, in addition: an AMS \texttt{amsart} preamble that restates the article's own declarations and macros used by the retained lines, namely the environments \texttt{proposition} on separate counters, together with the standard packages those lines need.  The retained environments print in this document as Proposition 1 in order of appearance, whereas the article prints the same items with the numbering of its own full document.  The complete native source of the article as distributed in the arXiv e-print for this exact version is bundled unchanged as \texttt{sources/191101/source0.tex}.
\begin{proposition}\label{section 3 main prop}
Let $\ell$ be an odd prime and let $n,r\geq 1$ be integers with $n\ge 2r$. Setting
$R=\Z/\ell^n\Z$, for $M=\begin{pmatrix}a&b\\ c&d\end{pmatrix}\in M_2(R)$ define
\[F(a,b,c,d):=(\tr M)^2-4\det M=(a-d)^2+4bc.
\]
Letting
\[
S_{n,r}:=\{M\in\op{GL}_2(R):\,v_\ell(F(a,b,c,d))\ge 2r\},
\]
we have that
\begin{equation}\label{S n r density}
\frac{|S_{n,r}|}{|\mathrm{GL}_2(\mathbb{Z}/\ell^n)|} 
= 2\ell^{-2r}
+
\frac{\ell^{-2r}\big(2+2\ell-(2r+2)\ell^{1-r}\big)}{\ell^{2}-1}.
\end{equation}
where the implied constant in $O_r$ depends on $r$ but not on $\ell$.
\end{proposition}

\begin{proof}
Make the invertible $R$–linear change of variables
\[
x=a-d,\qquad y=b,\qquad z=c,\qquad w=a+d,
\]
whose Jacobian is a unit in $R$ since $\ell$ is odd. Then with respect to the variables $(x,y,z,w)$, we find that
\[
F(a,b,c,d)=x^2+4yz,
\qquad\text{and}\quad
\det M=\frac{w^2-x^2}{4}-yz.
\]
We write $F(x,y,z):=F(a,b,c,d)=x^2+4yz$. The condition $v_\ell\left(F(a,b,c,d)\right)\ge2r$ defining the set $S_{n,r}$ depends only on $(x,y,z)$ and not on $w$. The congruence $x^2+4yz\equiv0\pmod{\ell^{2r}}$ implies that $\det M\equiv w^2/4\pmod{\ell^{2r}}$ and thus $\det M\in R^\times$ if and only if $w\not\equiv0\pmod\ell$.  Hence each
$(x,y,z)\in R^3$ satisfying $F(x,y,z)\equiv 0\pmod{\ell^{2r}}$ contributes to exactly $\ell^{\,n-1}(\ell-1)$ choices of $w$. Therefore, we find that
\[|S_{n,r}|=\ell^{\,n-1}(\ell-1)\#\{(x,y,z)\in R^3\mid F(x,y,z)\equiv 0\pmod{\ell^{2r}}\}.\]
\noindent The congruence $F(x,y,z)\equiv0\pmod{\ell^{2r}}$ depends only on residues
modulo $\ell^{2r}$, so the number of solutions modulo $\ell^n$ is
$\ell^{3(n-2r)}$ times the number
\[
N_0:=\#\{(x,y,z)\in(\Z/\ell^{2r}\Z)^3|\ F(x,y,z)\equiv0\pmod{\ell^{2r}}\}.
\] 
\noindent In other words,
\[|S_{n,r}|=\ell^{\,n-1}(\ell-1)\ell^{3(n-2r)} N_0=\ell^{4n-6r-1}(\ell-1)N_0.\]
\noindent Write $t=-x^2/4$ and $u=v_\ell(x)\in [0,2r]$. Thus
\begin{equation}\label{N_0 eqn}N_0=\sum_{u=0}^{2r} \#\left\{(x,y,z)\in(\Z/\ell^{2r}\Z)^3|\, v_\ell(x)=u\,\text{and}\, yz\equiv -\frac{x^2}{4}\pmod{\ell^{2r}}\right\}.\end{equation}
Suppose that $u<r$, the congruence 
\[yz\equiv -\frac{x^2}{4}\pmod{\ell^{2r}}\]determines $x^2\in \Z/\ell^{2r}\Z$. Writing $x=\ell^ux'$ where $x'$ is an element in $\Z/\ell^{2r-u}$ which is not divisible by $\ell$, we find that $(x')^2\pmod{\ell^{2r-2u}}$ is determined. Thus for any pair $(y,z)$ with $v_\ell(yz)=u$, the number of choices of $x$ is $2\ell^u$. Thus, we find that for $u<r$,
\begin{equation}\label{u<r eqn}\begin{split}&\#\left\{(x,y,z)\in(\Z/\ell^{2r}\Z)^3|\, v_\ell(x)=u\,\text{and}\, yz\equiv -\frac{x^2}{4}\pmod{\ell^{2r}}\right\}\\
=&2\ell^u\#\left\{(y,z)\in(\Z/\ell^{2r}\Z)^2|\, v_\ell(y)+v_\ell(z)=2u\right\}\\
=& 2\ell^u\sum_{i+j=2u} \#\left\{(y,z)\in(\Z/\ell^{2r}\Z)^2|\, v_\ell(y)=i\quad\text{and}\quad v_\ell(z)=j\right\}\\
=& 2\ell^u\sum_{i+j=2u} \varphi(\ell^{2r-i})\varphi(\ell^{2r-j})\\
=& 2\ell^u\sum_{i+j=2u} \ell^{4r-2u}\left(1-\ell^{-1}\right)^2\\
=& 2(2u+1) \ell^{4r-u}\left(1-\ell^{-1}\right)^2.
\end{split}\end{equation}
Next, consider the case when $u\geq r$ and thus $x^2\equiv 0\pmod{\ell^{2r}}$. In this case, we find that 
\[\begin{split}&\#\left\{(x,y,z)\in(\Z/\ell^{2r}\Z)^3|\, v_\ell(x)=u\,\text{and}\, yz\equiv 0\pmod{\ell^{2r}}\right\}\\
=&\sum_{\substack{0\leq i,j\leq 2r\\
i+j\geq 2r}} \#\left\{(x,y,z)\in(\Z/\ell^{2r}\Z)^2|v_\ell(x)=u,\quad \, v_\ell(y)=i\quad\text{and}\quad v_\ell(z)=j\right\}\\
=& \sum_{\substack{0\leq i,j\leq 2r\\
i+j\geq 2r}}\varphi(\ell^{2r-i})\varphi(\ell^{2r-j})\varphi(\ell^{2r-u}).\\
\end{split}\]
\noindent Changing variables \(p:=2r-i,\; q:=2r-j\), we find that \(p,q\in\{0,\dots,2r\}\) and the condition \(i+j\ge2r\) is equivalent to \(p+q\le2r\). Hence
\[
\sum_{\substack{0\le i,j\le 2r\\ i+j\ge 2r}}
\varphi(\ell^{2r-i})\varphi(\ell^{2r-j})
=
\sum_{p=0}^{2r}\sum_{q=0}^{2r-p}\varphi(\ell^{p})\varphi(\ell^{q}).
\]
\noindent For \(m\ge1\) we have \(\varphi(\ell^{m})=\ell^{m}-\ell^{m-1}=\ell^{m-1}(\ell-1)\), and \(\varphi(1)=1\). Set
\(
a_m:=\varphi(\ell^m)
\)
and denote the partial sums \(S_m:=\sum_{q=0}^{m}a_q\). Clearly, $S_0=1=\ell^0$ and for $m\geq 1$, one has that
\[
S_m=1+\sum_{q=1}^m \ell^{q-1}(\ell-1)=1+(\ell-1)\frac{\ell^{m}-1}{\ell-1}=\ell^{m}.
\]
Consequently,
\[
\begin{split} \sum_{p=0}^{2r}\sum_{q=0}^{2r-p}a_p a_q
=&\sum_{p=0}^{2r} a_p\Big(\sum_{q=0}^{2r-p}a_q\Big)
=\sum_{p=0}^{2r} a_p\,\ell^{2r-p}\\
=&\ell^{2r} + \sum_{p=1}^{2r} \ell^{p-1}(\ell-1)\,\ell^{2r-p}\\
= & \ell^{2r} + \sum_{p=1}^{2r} \ell^{2r-1}(\ell-1)\\
=&\ell^{2r-1}\big((2r+1)\ell - 2r\big).
\end{split}
\]
We have shown that for $u\geq r$,
\begin{equation}\label{u>=r eqn}\begin{split}
&\#\left\{(x,y,z)\in(\Z/\ell^{2r}\Z)^3|\, v_\ell(x)
=u\,\text{and}\, yz\equiv 0\pmod{\ell^{2r}}\right\}\\
=&\begin{cases}
\ell^{2r-1}\big((2r+1)\ell-2r\big), & \text{if }u=2r,\\
\ell^{4r-u-2}(\ell-1)\big((2r+1)\ell-2r\big), & \text{if }r\le u\le 2r-1.
\end{cases}
\end{split}
\end{equation}
Substituting \eqref{u<r eqn} and \eqref{u>=r eqn} into \eqref{N_0 eqn}, we find that
\[N_0=\sum_{u=0}^{r-1}2(2u+1) \ell^{4r-u}\left(1-\ell^{-1}\right)^2+\sum_{u=r}^{2r-1}\ell^{4r-u-2}(\ell-1)\big((2r+1)\ell-2r\big)+\ell^{2r-1}\big((2r+1)\ell-2r\big).\]
\noindent Set 
\[\begin{split}
&S_1:=\sum_{u=0}^{r-1}2(2u+1) \ell^{4r-u}\left(1-\ell^{-1}\right)^2,\\
&S_2:=\sum_{u=r}^{2r-1}\ell^{4r-u-2}(\ell-1)\big((2r+1)\ell-2r\big),\\
&S_3:=\ell^{2r-1}\big((2r+1)\ell-2r\big).
\end{split}\]
For the first sum above, consider the following reductions:
\[
\begin{split}\sum_{u=0}^{r-1} (2u+1)\ell^{-u}(1-\ell^{-1})^{2}
= &(1-\ell^{-1})^{2}\Bigl(2\sum_{u=0}^{r-1}u\ell^{-u}+\sum_{u=0}^{r-1}\ell^{-u}\Bigr)\\
= &(1-\ell^{-1})^{2}\left( 
2\cdot\frac{\ell^{-1}(1-r\ell^{-r-1}+(r-1)\ell^{-r})}{(1-\ell^{-1})^{2}}
+ \frac{1-\ell^{-r}}{1-\ell^{-1}}
\right)\\
= &2\ell^{-1}(1-r\ell^{-r-1}+(r-1)\ell^{-r}) + (1-\ell^{-r})(1-\ell^{-1})\\
= &1+\ell^{-1}-(2r+1)\ell^{-r}+(2r-1)\ell^{-r-1}.
\end{split}
\]
\noindent Multiply by $2 \ell^{4r}$ to deduce that
\[
S_1=\sum_{u=0}^{r-1}2(2u+1) \ell^{4r-u}\left(1-\ell^{-1}\right)^2= 2\ell^{4r} + 2\ell^{4r-1} - 2(2r+1)\ell^{3r} + 2(2r-1)\ell^{3r-1}.
\]
\noindent For the second sum, set $A = (2r+1)\ell - 2r$. We find that
\[
S_2=\sum_{u=r}^{2r-1} \ell^{4r-u-2} (\ell-1) A 
= A (\ell-1) \sum_{u=r}^{2r-1} \ell^{4r-u-2}=A \ell^{3r-1} - A \ell^{2r-1}.
\]
Thus one finds that
\[
S_2 + S_3 = A \ell^{3r-1} - A \ell^{2r-1} + A \ell^{2r-1} = A \ell^{3r-1}.
\]
Combining the above, we deduce that
\[
N_0= S_1 + S_2 + S_3 =  2\ell^{4r} + 2\ell^{4r-1} - (2r+1)\ell^{3r} - \ell^{3r-1}.
\]
and 
\[
|S_{n,r}| = \ell^{4n-6r-1}(\ell-1)N_0=(\ell-1)\, \ell^{4n-3r-2} \Big( 2 \ell^{r+1} + 2 \ell^r - (2r+1)\ell -1 \Big)
\]
Thus, we have shown that
\[\begin{split}
\frac{|S_{n,r}|}{|\op{GL}_2(\mathbb{Z}/\ell^n)|} 
=&\frac{(\ell-1)\,\ell^{4n-3r-2}\big(2\ell^{r+1}+2\ell^{r}-(2r+1)\ell-1\big)}
       {\ell^{4n-3}(\ell-1)^2(\ell+1)}\\
=&\frac{\ell^{\,1-3r}\big(2\ell^{r+1}+2\ell^{r}-(2r+1)\ell-1\big)}
       {(\ell^2-1)}\\
=&
2\ell^{-2r}
+
\frac{\ell^{-2r}\big(2+2\ell-(2r+2)\ell^{1-r}\big)}{\ell^{2}-1}.
\end{split}
\]
\end{proof}

\end{document}
